\stagechapter{C}{Codebook and pilot}

\noindent\textbf{What this stage covers.} How the coding instrument is built and tested before the main study: choice of analytic family, origin of the codes, structure of a codebook entry, unit of analysis, pilot rounds and the rule for freezing. The draft thesis plan codes units of one paragraph or one coherent block of interface text on three axes (A coverage, B stance, C protection), defines each label with a fitting example and a counter-example, runs a five-app pilot (three fitness, two fertility) and freezes the release before the main study. Related decisions: D8 (two-coder procedure), D10 (checklist and outputs).

\section{Choosing the analytic family}
\pa{braun2021} separates coding-reliability, codebook and reflexive thematic analysis and assigns quality criteria to each; a predefined codebook with independent coders and reliability statistics belongs to the first two families, and \pa{braun2006} should not be cited as authority for it. \pa{hsieh2005} offers the vocabulary: the thesis codebook is \emph{directed} (codes from prior work), possibly with a \emph{summative} part (counting terms such as ``may''). \pa{mayring2000} is the German-tradition name for the same procedure (deductive category application with a coding agenda: category, definition, anchor example, coding rule, plus a revision loop); \pa{krippendorff2019} is the textbook on designs, units and the codebook as a recording instruction. \pa{fereday2006} supports a hybrid: a priori template codes plus documented inductive additions. \pa{mcdonald2019} and \pa{oconnor2020} explain when reliability statistics are expected: a predetermined codebook applied to a defined set of observations.

\section{Where the codes come from}
\begin{itemize}
\item \emph{Policy annotation schemes.} \pa{wilson2016} built OPP-115 with expert annotators, categories plus category-specific attributes and an explicit ``unspecified'' value so vague policies are coded as vague; \pa{harkous2018} reuses those categories for automatic annotation.
\item \emph{Privacy-policy appraisal instruments.} \pa{lamonica2021} (see the form \pp{lamonica2021tool}), \pa{huckvale2015} (policy schema in \pp{huckvale2015af}), \pa{hutton2018}, \pa{apa2024} and \pa{benjumea2020} (a scoping review of assessment criteria, with the finding that decision rules and the place where each item is looked for are the main reproducibility gap).
\item \emph{Study-specific codebooks.} \pa{malki2024} lists code categories with a definition and a real quotation per code (Table C.1 of its paper); \pa{shipp2020} defines hedging language operationally with word lists and sentence-in-context reading; \pa{zou2019} codes risk language (hedges, ``no evidence'' claims, links to consequences) in 161 breach notices; \pa{doan2025} codes contextual-integrity parameters and flags flows as missing, bloated or vague with a published vague-term list.
\item \emph{Theory-based dimensions.} \pa{nissenbaum2004} and \pa{martin2023} for flow parameters and purposes; \pa{kelley2009} for the data--purpose--recipient triplet; \pa{schaub2015} for timing, channel, modality and control of notices; \pa{wp292018} for clarity, accessibility and content criteria (annex rows for recipients, retention, rights, consent withdrawal).
\item \emph{Interface codes.} \pa{mathur2019} gives five characteristics of dark patterns (asymmetric, covert, deceptive, information hiding, restrictive) that describe an element without claiming intent; \pa{gray2018} gives five strategies (nagging, obstruction, sneaking, interface interference, forced action), built by applying a typology deductively and then recoding openly; \pa{nouwens2020} derives countable consent conditions from law.
\end{itemize}

\section{Structure of a codebook entry}
\pa{macqueen1998} gives a six-part entry (code, brief definition, full definition, when to use, when not to use, examples) and a refinement loop in which disagreements are triaged into coder error and definition problems; the exclusion criteria and cross-references to neighbouring codes target confusion between similar constructs. \pa{fereday2006} suggests label, definition and ``how to know it occurs'' as a minimal form. \pa{hutton2018} and \pa{stoyanov2015} show \emph{anchored} levels for ordinal ratings (an observable state per level) and an N/A rule that adjusts the denominator. Additions suggested by the cards: a field for observable evidence and screenshot reference (walkthrough codes), the data type and recipient where applicable, and the language of the source. The draft thesis entry (definition, example, counter-example) matches the MacQueen core.

\section{Units and unitizing}
\pa{krippendorff2019} asks to separate sampling, recording and context units and to measure unitizing reliability if units are judged; \pa{campbell2013} proposes that the principal investigator brackets the units first, then raises discriminant capability through independent coding and negotiated agreement (reported separately from independent reliability). The walkthrough steps are pre-segmented by the protocol, so unitizing is mainly a policy-text problem: the draft's ``paragraph or coherent block'' with context available, and supporting quotations identified precisely, follows this logic.

\section{The pilot}
\begin{itemize}
\item \pa{hruschka2004} shows empirically that agreement on a first draft is poor and rises over rounds (two to nine rounds in its three CDC studies, fewer when each round had under 20 codes); plan several rounds on fresh subsets, keep the number of codes per pass below about 20 (split into families), and state the stop rule and cut-off beforehand.
\item \pa{lombard2002} suggests a pilot of at least 30 units for the reliability check before the main coding; \pa{zou2019} piloted with three coders on a random 12\,\% subset and reported Fleiss' $\kappa=.75$ before splitting the corpus; \pa{malki2024} piloted the deductive draft on five random policies, then iterated with independent coding of 3--5 policies until saturation, but calculated no inter-rater reliability.
\item \pa{stoyanov2015} is the template for instrument development: derive items from literature criteria, expert-panel classification, anchors, pilot, documented revisions, then a new sample (its pilot ICC of .57 shows why a revision log matters).
\item \pa{nowell2017} supplies documentation: audit trail, versioned codebook with change notes, coder-discussion minutes.
\end{itemize}

\begin{table}[t]
\centering\small
\caption{Codebook design choices in the corpus.}
\label{tab:C-design}
\begin{tabularx}{\textwidth}{@{}P{2.5cm}P{3.3cm}YY@{}}
\toprule
\textbf{Source} & \textbf{Origin of codes} & \textbf{Entry and unit} & \textbf{Pilot and reliability}\\
\midrule
\pa{malki2024} & deductive draft (FemTech work, GDPR rights) & definition plus real quote per code; line-by-line in MAXQDA & 5-policy pilot; no IRR\\
\pa{wilson2016} & expert design; categories and attributes & segment-level, ``unspecified'' value & three annotators, Fleiss $\kappa$ per category\\
\pa{hutton2018} & 26 heuristics from the literature & anchored levels, N/A rule & second rater on a random subset\\
\pa{zou2019} & breach-notice features & hedges, reassurance, consequences & 3 coders, 12\,\% subset, $\kappa=.75$\\
\pa{doan2025} & CI parameters plus bottom-up codes & flows flagged missing, bloated, vague & one coder with consultation (gap)\\
\pa{shipp2020} & scope categories (Sunyaev) & word-list hedging definition & see card\\
\pa{macqueen1998} & team-based refinement & six-part entry & independent coding of same sample\\
\bottomrule
\end{tabularx}
\end{table}

\begin{thesisbox}[Implications for this thesis (stage C)]
\begin{enumerate}[leftmargin=1.3em]
\item Name the approach as a directed, codebook-based content analysis with a coding agenda (\pa{mayring2000}, \pa{krippendorff2019}) and say that it is not reflexive thematic analysis (\pa{braun2021}); record the choice in the Decision Log.
\item Use the MacQueen entry with two thesis-specific fields (observable evidence, source language) and anchored levels for Axis C strengths; keep stance codes linked to the specific harm and supporting span, as the draft says.
\item Seed and cross-check the codes against OPP-115, LaMonica, Huckvale AF4, WP260 annex and the dark-pattern codes, and document any code that has no precedent.
\item Plan the pilot as several rounds (budget for more than two) with a stop rule written before it starts, a revision log and a versioned freeze; keep pilot apps in the sample but out of the validation set (as the draft states).
\end{enumerate}
\end{thesisbox}
